\chapter{Funciones de referencia}\label{ch:g10:reffunc}

Un puñado de funciones sencillas (afín, cuadrado, inversa, raíz cuadrada,
cubo) aparecen por todas partes, solas o combinadas. Saberse de memoria sus
\omterm{def:g10:functions:graph}{gráficas} y su variación convierte muchos problemas en un esbozo rápido.
Este capítulo estudia cada una de ellas y demuestra su variación con la
técnica de comparación del \cref{ch:g10:functions}.

\section{Funciones afines}

\begin{definition}[Función afín]\label{def:g10:reffunc:affine}
Una \emph{función afín}\index{función afín} es una \omterm{def:g10:functions:function}{función} de la forma
\[
f(x) = mx + p,
\]
donde $m$ y $p$ son \omterm{def:g10:numbers:sets}{números reales} fijos. Su \omterm{def:g10:functions:graph}{gráfica} es una recta: $m$ es
la \emph{pendiente}\index{pendiente} y $p$, la \emph{ordenada en el
origen}\index{ordenada en el origen} (la recta corta al eje vertical en
$(0, p)$). Cuando $p = 0$, la \omterm{def:g10:functions:function}{función} es \emph{lineal}: $f(x) = mx$
expresa una proporcionalidad.
\end{definition}

\begin{proposition}[Pendiente y variación]\label{prop:g10:reffunc:affine}
Sea $f(x) = mx + p$.
\begin{enumerate}
  \item Para dos entradas distintas cualesquiera $u \neq v$:
        $m = \dfrac{f(v) - f(u)}{v - u}$ (la \omterm{def:g10:reffunc:affine}{pendiente} es la variación del
        resultado por unidad de variación de la entrada).
  \item Si $m > 0$, $f$ es estrictamente \omterm{def:g10:functions:variations}{creciente} en $\R$; si $m < 0$,
        estrictamente \omterm{def:g10:functions:variations}{decreciente}; si $m = 0$, constante.
\end{enumerate}
\end{proposition}

\begin{proof}
1. Calculamos $f(v) - f(u) = (mv + p) - (mu + p) = m(v - u)$ y dividimos
entre $v - u \neq 0$.

2. Si $u < v$, entonces $f(v) - f(u) = m(v-u)$ tiene el signo de $m$,
puesto que $v - u > 0$: con $m$ positivo, $f(u) < f(v)$ (\omterm{def:g10:functions:variations}{creciente}); con
$m$ negativo, $f(u) > f(v)$ (\omterm{def:g10:functions:variations}{decreciente}).
\end{proof}

\begin{omfigure}
\begin{tikzpicture}
\begin{axis}[omaxis,
  width=7.8cm, height=5.6cm,
  xmin=-3.4, xmax=3.4, ymin=-2.8, ymax=4.4,
  xtick={-2,-1,1,2}, ytick={-1,1,2,3}]
  \addplot[omDef, thick, domain=-3:2.6] {0.5*x + 1};
  \addplot[omThm, thick, domain=-2.2:3] {-x + 1};
  \node[omDef, font=\small] at (axis cs:2.5,2.7) {$y = \tfrac12 x + 1$};
  \node[omThm, font=\small] at (axis cs:-2.0,2.6) {$y = -x + 1$};
  \draw[dashed] (axis cs:1,1.5) -- (axis cs:2,1.5) -- (axis cs:2,2);
  \node[font=\footnotesize, below] at (axis cs:1.5,1.5) {$1$};
  \node[font=\footnotesize, right] at (axis cs:2,1.75) {$m$};
\end{axis}
\end{tikzpicture}

{\small Dos funciones afines: \omterm{def:g10:functions:variations}{creciente} para $m = \frac12 > 0$ y
\omterm{def:g10:functions:variations}{decreciente} para $m = -1 < 0$. Al avanzar una unidad hacia la derecha, $y$
varía en $m$.}
\end{omfigure}

\begin{example}\label{ex:g10:reffunc:affine}
Hallemos la \omterm{def:g10:reffunc:affine}{función afín} cuya \omterm{def:g10:functions:graph}{gráfica} pasa por $A(1, 3)$ y $B(4, 9)$. La
\omterm{def:g10:reffunc:affine}{pendiente} es
\[
m = \frac{9 - 3}{4 - 1} = \frac{6}{3} = 2 .
\]
Entonces $f(x) = 2x + p$, y $f(1) = 3$ da $2 + p = 3$, luego $p = 1$:
$f(x) = 2x + 1$. Comprobación con $B$: $f(4) = 9$.
\end{example}

\section{La función cuadrado}

\begin{proposition}[Función cuadrado]\label{prop:g10:reffunc:square}
La \omterm{def:g10:functions:function}{función} $f(x) = x^2$, definida en $\R$:
\begin{enumerate}
  \item es estrictamente \omterm{def:g10:functions:variations}{decreciente} en $\intoc{-\infty}{0}$ y
        estrictamente \omterm{def:g10:functions:variations}{creciente} en $\intco{0}{+\infty}$, con \omterm{def:g10:functions:extrema}{mínimo} $0$ en
        $x = 0$;
  \item cumple $f(-x) = f(x)$ para todo $x$: su \omterm{def:g10:functions:graph}{gráfica}, una
        \emph{parábola}\index{parábola}, es simétrica respecto del eje
        vertical.
\end{enumerate}
\end{proposition}

\begin{proof}
1. Sea $0 \leq u < v$. Entonces
\[
f(v) - f(u) = v^2 - u^2 = (v - u)(v + u) > 0,
\]
puesto que $v - u > 0$ y $v + u > 0$: $f$ es estrictamente \omterm{def:g10:functions:variations}{creciente} en
$\intco{0}{+\infty}$. Si $u < v \leq 0$, entonces $v - u > 0$ pero
$v + u < 0$, luego $f(v) - f(u) < 0$: estrictamente \omterm{def:g10:functions:variations}{decreciente}. Por
último, $x^2 \geq 0 = f(0)$ para todo $x$.

2. $(-x)^2 = x^2$; así que los puntos $(x, x^2)$ y $(-x, x^2)$, simétricos
respecto del eje vertical, están los dos en la \omterm{def:g10:functions:graph}{gráfica}.
\end{proof}

\begin{remark}[Cuadrados y orden]
Como la \omterm{def:g10:functions:function}{función} cuadrado decrece en el lado negativo, elevar al cuadrado
\emph{invierte} el orden de los números negativos: $-3 < -2$, pero
$9 > 4$. No eleves nunca al cuadrado los dos miembros de una desigualdad
sin comprobar antes los signos.
\end{remark}

\section{La función inversa}

\begin{proposition}[Función inversa]\label{prop:g10:reffunc:inverse}
La \omterm{def:g10:functions:function}{función} $f(x) = \dfrac1x$, definida para $x \neq 0$:
\begin{enumerate}
  \item es estrictamente \omterm{def:g10:functions:variations}{decreciente} en $\intoo{-\infty}{0}$ y
        estrictamente \omterm{def:g10:functions:variations}{decreciente} en $\intoo{0}{+\infty}$;
  \item cumple $f(-x) = -f(x)$: su \omterm{def:g10:functions:graph}{gráfica}, una
        \emph{hipérbola}\index{hipérbola}, es simétrica respecto del
        origen.
\end{enumerate}
\end{proposition}

\begin{proof}
1. Sea $0 < u < v$. Entonces
\[
f(v) - f(u) = \frac1v - \frac1u = \frac{u - v}{uv} < 0,
\]
puesto que $u - v < 0$ y $uv > 0$: estrictamente \omterm{def:g10:functions:variations}{decreciente} en
$\intoo{0}{+\infty}$. En $\intoo{-\infty}{0}$, el mismo cociente tiene de
nuevo $u - v < 0$ y $uv > 0$ (producto de dos negativos), así que $f$
también decrece allí.

2. $\frac{1}{-x} = -\frac1x$.
\end{proof}

\begin{remark}
Cuidado: $\frac1x$ \emph{no} es \omterm{def:g10:functions:variations}{decreciente} en todo su \omterm{def:g10:functions:function}{dominio}. De
$u = -1$ a $v = 1$ el valor salta de $-1$ a $1$. Hay que estudiar las dos
ramas por separado.
\end{remark}

\begin{omfigure}
\begin{tikzpicture}
\begin{axis}[omaxis,
  width=5.6cm, height=5.2cm,
  xmin=-3.2, xmax=3.2, ymin=-2.6, ymax=6.4,
  xtick={-2,-1,1,2}, ytick={1,4},
  title={$y = x^2$}]
  \addplot[omDef, thick, domain=-2.45:2.45, samples=90] {x^2};
\end{axis}
\end{tikzpicture}\hfill
\begin{tikzpicture}
\begin{axis}[omaxis,
  width=5.6cm, height=5.2cm,
  xmin=-3.4, xmax=3.4, ymin=-3.4, ymax=3.4,
  xtick={-1,1}, ytick={-1,1},
  title={$y = \frac1x$}]
  \addplot[omDef, thick, domain=0.3:3.1, samples=80] {1/x};
  \addplot[omDef, thick, domain=-3.1:-0.3, samples=80] {1/x};
\end{axis}
\end{tikzpicture}

{\small La parábola $y = x^2$ (simétrica respecto del eje vertical) y la
hipérbola $y = \frac1x$ (dos ramas, simétricas respecto del origen).}
\end{omfigure}

\section{Raíz cuadrada y cubo}

\begin{proposition}[Función raíz cuadrada]\label{prop:g10:reffunc:sqrt}
La \omterm{def:g10:functions:function}{función} $f(x) = \sqrt{x}$, definida en $\intco{0}{+\infty}$, es
estrictamente \omterm{def:g10:functions:variations}{creciente}.
\end{proposition}

\begin{proof}
Sea $0 \leq u < v$. Multiplicamos y dividimos por el \emph{conjugado}:
\[
\sqrt v - \sqrt u
= \frac{(\sqrt v - \sqrt u)(\sqrt v + \sqrt u)}{\sqrt v + \sqrt u}
= \frac{v - u}{\sqrt v + \sqrt u} > 0,
\]
puesto que $v - u > 0$ y $\sqrt v + \sqrt u > 0$ (obsérvese que $v > 0$,
luego $\sqrt v > 0$).
\end{proof}

\begin{proposition}[Función cubo]\label{prop:g10:reffunc:cube}
La \omterm{def:g10:functions:function}{función} $f(x) = x^3$, definida en $\R$, es estrictamente \omterm{def:g10:functions:variations}{creciente} y su
\omterm{def:g10:functions:graph}{gráfica} es simétrica respecto del origen.
\end{proposition}
\admitted

\begin{omfigure}
\begin{tikzpicture}
\begin{axis}[omaxis,
  width=5.6cm, height=5.0cm,
  xmin=-0.7, xmax=5.4, ymin=-0.7, ymax=2.9,
  xtick={1,4}, ytick={1,2},
  title={$y = \sqrt x$}]
  \addplot[omDef, thick, domain=0:5.1, samples=100] {sqrt(x)};
\end{axis}
\end{tikzpicture}\hfill
\begin{tikzpicture}
\begin{axis}[omaxis,
  width=5.6cm, height=5.0cm,
  xmin=-2.2, xmax=2.2, ymin=-3.4, ymax=3.4,
  xtick={-1,1}, ytick={-1,1},
  title={$y = x^3$}]
  \addplot[omDef, thick, domain=-1.48:1.48, samples=90] {x^3};
\end{axis}
\end{tikzpicture}

{\small La raíz cuadrada (definida solo para $x \geq 0$) y el cubo, las dos
\omterm{def:g10:functions:variations}{crecientes}.}
\end{omfigure}

\begin{proposition}[Comparación de $x$, $x^2$ y $\sqrt x$]\label{prop:g10:reffunc:compare}
Para $0 < x < 1$:\ \ $x^2 < x < \sqrt x$.\quad
Para $x > 1$:\ \ $\sqrt x < x < x^2$.\quad
En $x = 0$ y $x = 1$ los tres valores coinciden.
\end{proposition}

\begin{proof}
Supongamos $0 < x < 1$. Multiplicando $x < 1$ por $x > 0$ se obtiene
$x^2 < x$. Después, $x < \sqrt x$: como los dos miembros son positivos y
elevar al cuadrado es \omterm{def:g10:functions:variations}{creciente} sobre los positivos, esta desigualdad
equivale a $x^2 < x$, que acabamos de demostrar. Para $x > 1$,
multiplicando $x > 1$ por $x$ se obtiene $x^2 > x$, y $\sqrt x < x$ se
sigue con el mismo argumento del cuadrado.
\end{proof}

\begin{omfigure}
\begin{tikzpicture}
\begin{axis}[omaxis,
  width=7.4cm, height=5.8cm,
  xmin=-0.35, xmax=2.5, ymin=-0.35, ymax=2.9,
  xtick={1,2}, ytick={1,2}]
  \addplot[omThm, thick, domain=0:1.65, samples=80] {x^2};
  \addplot[omDef, thick, domain=0:2.35] {x};
  \addplot[omMeth, thick, domain=0:2.35, samples=80] {sqrt(x)};
  \fill (axis cs:1,1) circle (1.5pt);
  \node[omThm, font=\small, right] at (axis cs:1.45,2.4) {$y = x^2$};
  \node[omDef, font=\small, right] at (axis cs:2.0,1.85) {$y = x$};
  \node[omMeth, font=\small, right] at (axis cs:2.0,1.28) {$y = \sqrt x$};
\end{axis}
\end{tikzpicture}

{\small Las tres curvas se cortan en $(0,0)$ y $(1,1)$ e intercambian allí
su orden: entre $0$ y $1$, $x^2$ es la menor y $\sqrt x$ la mayor; a
partir de $1$ el orden se invierte.}
\end{omfigure}

\begin{method}[Comparar imágenes]\label{met:g10:reffunc:compare}
Para comparar $f(a)$ y $f(b)$ con una \omterm{def:g10:functions:function}{función} de referencia $f$:
\begin{enumerate}
  \item sitúa $a$ y $b$ en un \omterm{def:g10:numbers:interval}{intervalo} donde se conozca la variación de
        $f$;
  \item si $f$ crece allí, las imágenes quedan en el mismo orden que las
        entradas; si $f$ decrece, el orden se invierte;
  \item si $a$ y $b$ no están en el mismo \omterm{def:g10:numbers:interval}{intervalo} de monotonía, ocúpate
        primero de los signos (por ejemplo, para el cuadrado: una entrada
        negativa y otra positiva).
\end{enumerate}
\end{method}

\begin{example}
Comparemos $\dfrac{1}{2.3}$ y $\dfrac{1}{2.4}$: las dos entradas están en
$\intoo{0}{+\infty}$, donde la \omterm{def:g10:functions:function}{función} inversa decrece, y $2.3 < 2.4$,
luego $\dfrac{1}{2.3} > \dfrac{1}{2.4}$. Comparemos $(-1.2)^2$ y
$(-1.5)^2$: en $\intoc{-\infty}{0}$ el cuadrado decrece y
$-1.5 < -1.2$, luego $(-1.5)^2 > (-1.2)^2$; en efecto, $2.25 > 1.44$.
\end{example}

\section{Ejercicios}

\begin{exercise}[$\star$]\label{exo:g10:reffunc:1}
Para cada \omterm{def:g10:reffunc:affine}{función afín}, da su \omterm{def:g10:reffunc:affine}{pendiente} y su \omterm{def:g10:reffunc:affine}{ordenada en el origen} y di si
es \omterm{def:g10:functions:variations}{creciente} o \omterm{def:g10:functions:variations}{decreciente}:
\[
f(x) = 3x - 2, \qquad
g(x) = -\tfrac12 x + 4, \qquad
h(x) = 7, \qquad
k(x) = -x .
\]
\end{exercise}

\begin{exercise}[$\star$]\label{exo:g10:reffunc:2}
Halla la \omterm{def:g10:reffunc:affine}{función afín} cuya \omterm{def:g10:functions:graph}{gráfica} pasa por $(2, 1)$ y $(5, 10)$, y
después la que pasa por $(-1, 4)$ y $(3, -4)$.
\end{exercise}

\begin{exercise}[$\star$]\label{exo:g10:reffunc:3}
Sin calculadora, compara:
\[
(3.1)^2 \text{ y } (3.2)^2; \qquad
(-2.7)^2 \text{ y } (-2.8)^2; \qquad
\frac{1}{5.1} \text{ y } \frac{1}{5.2}; \qquad
\sqrt{17} \text{ y } \sqrt{15}.
\]
\end{exercise}

\begin{exercise}[$\star$]\label{exo:g10:reffunc:4}
Usando la \omterm{def:g10:functions:graph}{gráfica} de la \omterm{def:g10:functions:function}{función} cuadrado, resuelve $x^2 = 16$, después
$x^2 \leq 16$ y después $x^2 > 9$.
\end{exercise}

\begin{exercise}[$\star$]\label{exo:g10:reffunc:5}
Sea $x \in \intoo{0}{1}$. Ordena los números $x$, $x^2$, $x^3$ y $\sqrt x$
de menor a mayor y comprueba tu respuesta con $x = 0.25$.
\end{exercise}

\begin{exercise}[$\star\star$]\label{exo:g10:reffunc:6}
Resuelve gráficamente y después algebraicamente: $\dfrac1x = 2$ y
$\dfrac1x < 2$ para $x > 0$. ¿Qué cambia si se admite también $x < 0$?
\end{exercise}

\begin{exercise}[$\star\star$]\label{exo:g10:reffunc:7}
Sabiendo que la \omterm{def:g10:functions:function}{función} cuadrado decrece en $\intoc{-\infty}{0}$ y crece
en $\intco{0}{+\infty}$, acota $x^2$ cuando:
\[
\text{(a) } x \in \intcc{2}{5};
\qquad
\text{(b) } x \in \intcc{-3}{-1};
\qquad
\text{(c) } x \in \intcc{-2}{3}.
\]
(En el caso (c), atención: el \omterm{def:g10:functions:extrema}{mínimo} de $x^2$ no está en un extremo.)
\end{exercise}

\begin{exercise}[$\star\star$]\label{exo:g10:reffunc:8}
Una compañía de taxis cobra una cuota fija de $4$ más $1.5$ por
kilómetro; otra no cobra cuota y pide $2$ por kilómetro. Modeliza cada
precio con una \omterm{def:g10:reffunc:affine}{función afín} de la distancia $x$, representa las dos rectas
y halla a partir de qué distancia sale más barata la primera compañía.
\end{exercise}

\begin{exercise}[$\star\star$]\label{exo:g10:reffunc:9}
Demuestra que para todos $a, b \geq 0$ se cumple
$\sqrt{a + b} \leq \sqrt a + \sqrt b$. (Indicación: los dos miembros son
no negativos, así que compara sus cuadrados.) ¿Cuándo hay igualdad?
\end{exercise}

\begin{exercise}[$\star\star\star$]\label{exo:g10:reffunc:10}
Sea $f(x) = \dfrac{2x + 1}{x - 1}$, definida para $x \neq 1$.
\begin{enumerate}
  \item Demuestra que $f(x) = 2 + \dfrac{3}{x - 1}$ para todo $x \neq 1$.
  \item Deduce la variación de $f$ en $\intoo{1}{+\infty}$ a partir de la
        de la \omterm{def:g10:functions:function}{función} inversa.
\end{enumerate}
\end{exercise}

\section{Problema: Las leyes de la naturaleza hablan en funciones de referencia}

\begin{problem}[{Problema de fin de semana --- las distancias de frenado
crecen como $v^2$, las palancas obedecen a $1/d$, los planetas siguen
$a^{3/2}$: leer la física con las funciones de referencia}]
\label{pb:g10:reffunc:1}
Abre un libro de física y verás las mismas \omterm{def:g10:functions:graph}{gráficas} en todas las páginas:
la parábola, la hipérbola, la curva de la raíz. La naturaleza escribe sus
leyes con las funciones de referencia de este capítulo, y conocer sus
formas (\cref{prop:g10:reffunc:square}, \cref{prop:g10:reffunc:inverse},
\cref{prop:g10:reffunc:sqrt}, \cref{prop:g10:reffunc:cube}) basta para
frenar un coche, equilibrar una palanca y cronometrar los planetas. El
gran final es la tercera ley de Kepler, leída directamente en los datos
del sistema solar.

\medskip
\noindent\textbf{Parte I --- La ley cuadrática del frenado.}
Una regla práctica para la distancia de frenado de un coche sobre asfalto
seco: $d = \left(\frac{v}{10}\right)^{2}$ metros, a una velocidad de $v$
km/h.
\begin{enumerate}
  \item Calcula las distancias de frenado a $30$, $50$, $90$ y $130$ km/h.
  \item Al doblar la velocidad, ¿por cuánto queda multiplicada la
        distancia de frenado? Explícalo con el comportamiento de escala de
        la \omterm{def:g10:functions:function}{función} cuadrado y compruébalo con $50 \to 100$ km/h.
  \item Investigación de un accidente: las marcas de frenada miden $50$ m.
        ¿De qué velocidad acusa la fórmula al conductor (redondeando al
        km/h)? ¿Qué \omterm{def:g10:functions:function}{función} de referencia ha respondido, y en qué \omterm{def:g10:functions:function}{dominio}
        es inequívoca la lectura (\cref{prop:g10:reffunc:sqrt})?
  \item Compara la distancia adicional que causa \emph{un} km/h más a
        velocidad baja y a velocidad alta: calcula $d(31) - d(30)$ y
        $d(91) - d(90)$. ¿Qué propiedad de la parábola
        (\cref{prop:g10:reffunc:square}) ilustran las dos respuestas?
  \item En una frenada real hay que añadir el tiempo de reacción, un
        segundo aproximadamente, durante el cual el coche recorre $0.28v$
        metros: $D(v) = 0.28v + \left(\frac{v}{10}\right)^2$. Calcula
        $D(50)$ y $D(130)$. ¿Cuál de los dos términos, el afín o el
        cuadrático, domina a velocidad urbana? ¿Y en autopista?
\end{enumerate}

\medskip
\noindent\textbf{Parte II --- Las hipérbolas de la vida diaria.}
\begin{enumerate}[resume]
  \item Una palanca está en equilibrio cuando fuerza $\times$ distancia es
        igual a los dos lados. Un niño de $60$ kg se sienta a $2$ m del
        punto de apoyo; la fuerza que equilibra a distancia $d$ del otro
        lado es $F(d) = \frac{120}{d}$ (en equivalentes de kilogramo).
        Calcula $F(0.5)$, $F(1)$, $F(3)$.
  \item ¿Qué dice la variación de la \omterm{def:g10:functions:function}{función} inversa
        (\cref{prop:g10:reffunc:inverse}) sobre las palancas? ¿Y por qué
        se esconde en la cola de la hipérbola la bravata de Arquímedes
        (``dadme un punto de apoyo y moveré la Tierra'')? ¿Qué prohíbe
        $d = 0$?
  \item El sonido y la luz se atenúan con el \emph{cuadrado} de la
        distancia: a $1$ m de una lámpara la intensidad es de $100$
        unidades; dala a $2$, $3$ y $10$ m. Explica el exponente $2$ con
        una esfera: ¿sobre qué superficie se ha repartido la luz a
        distancia $d$ (problema de la lápida de Arquímedes, en el volumen
        anterior)?
  \item Un autobús para la excursión escolar cuesta $600$ euros, a
        repartir a partes iguales entre $n$ participantes. Haz una tabla
        del coste por cabeza para $n = 10, 20, 30$, halla cuántos
        participantes lo bajan a $25$ euros o menos y di qué promete, y
        qué se niega a dar, la forma de la hipérbola cuando $n$ crece.
  \item Producir $x$ carteles cuesta $2x + 3$ euros ($3$ euros de
        preparación). Demuestra que el coste \emph{medio} por cartel es
        $a(x) = 2 + \frac3x$, describe su variación e interpreta
        económicamente su asíntota horizontal. (Compara con el álgebra del
        \cref{exo:g10:reffunc:10}.)
\end{enumerate}

\medskip
\noindent\textbf{Parte III --- La armonía de Kepler.}
Para cada planeta, sea $a$ su distancia media al Sol (en unidades
astronómicas, Tierra $= 1$) y $T$ su periodo (en años). Datos: Marte
$a = 1.52$, $T = 1.88$; Júpiter $a = 5.20$, $T = 11.86$; Saturno
$a = 9.54$, $T = 29.4$.
\begin{enumerate}[resume]
  \item Calcula $a^3$ y $T^2$ para la Tierra, Marte y Júpiter. ¿Qué
        observó Kepler en 1618?
  \item Enuncia la ley como una \omterm{def:g10:functions:function}{función}: $T = \sqrt{a^3} = a\sqrt a$.
        Pruébala con Saturno.
  \item Dos predicciones: un asteroide orbita a $a = 4$ UA; halla su
        periodo (los números salen enteros). Un cometa vuelve cada $27$
        años; halla su distancia media (busca un cubo perfecto).
  \item ¿Qué tres funciones de referencia trenza esta ley? ¿Y cuál es su
        regla de escala: si $a$ se multiplica por $4$, qué le ocurre a
        $T$? (Compruébalo con el asteroide frente a la Tierra.)
  \item Kepler leyó esta ley en los datos de Tycho Brahe setenta años
        antes de que la gravitación de Newton la explicara. En una frase:
        ¿qué dice este episodio sobre el poder de reconocer una \omterm{def:g10:functions:function}{función} de
        referencia en una tabla de números?
\end{enumerate}

\medskip
\noindent\textbf{Parte IV --- La liebre y la tortuga.}
\begin{enumerate}[resume]
  \item Ordena $\sqrt x$, $x$ y $x^2$ en $\intoo{0}{1}$ y en
        $\intoo{1}{+\infty}$ (\cref{prop:g10:reffunc:compare}) y comprueba
        las dos ordenaciones en $x = 0.25$ y en $x = 4$.
  \item Resuelve $\sqrt x > x$ por completo (para $x \geq 0$, eleva al
        cuadrado con cuidado y termina con una \omterm{def:g10:algebra:signtable}{tabla de signos}).
  \item Entra el cubo: ordena $x$, $x^2$, $x^3$ en $x = 0.5$ y en $x = 2$
        y halla todos los puntos donde se cortan las funciones cuadrado y
        cubo.
  \item Dos planes de ahorro pagan, al cabo de $t$ años, $100\sqrt t$
        euros (plan A) o $10 t^2$ euros (plan B). Demuestra que B adelanta
        a A cuando $t^3 = 100$ y da el momento del cruce con precisión de
        un mes. ¿Moraleja sobre raíces frente a potencias a largo plazo?
  \item Final: para cada ley de este problema (el frenado con $v^2$, la
        palanca con $1/d$, la lámpara con $1/d^2$, Kepler con $a^{3/2}$,
        el plan A con $\sqrt t$), di en una frase qué rasgo de su \omterm{def:g10:functions:graph}{gráfica}
        de referencia lleva el significado físico (subida cada vez más
        pronunciada, asíntota vertical, esfera que se expande, exponente
        de escala, crecimiento que se aplana). Conclusión: las funciones
        de referencia son el alfabeto; la naturaleza escribe con ellas.
\end{enumerate}
\end{problem}
